%------------------------------------------------------------------------------%
% \chapter{Plano Cartesiano e Equações}
%------------------------------------------------------------------------------%

%------------------------------------------------------------------------------%
\begin{exercicios}{Distância Entre Pontos}
    
    \begin{exercicio}
        Calcule a distância entre os pontos:
        \begin{alphalist}[2]
            \item $(4,1)$ e $(2,3)$
            \item $(1,2)$ e $(-2,-4)$
            \item $(-5,-2)$ e $(3,-2)$
            \item $(-4,5)$ e $(1,-10)$
            \item $(6,4)$ e $(-3,1)$
            \item $(-3,4)$ e $(7,2)$
            \item $(-2,-6)$ e $(-1,1)$
            \item $(5,-2)$ e $(-9,5)$
        \end{alphalist}
        \begin{answers}
            \begin{mathlist}[4]
                \item 2 \sqrt{2}
                \item 3 \sqrt{5}
                \item 8
                \item 5 \sqrt{10}
                \item 3 \sqrt{10}
                \item 2 \sqrt{26}
                \item 5 \sqrt{2}
                \item 7 \sqrt{5}
            \end{mathlist}
        \end{answers}
    \end{exercicio}
    
    \begin{exercicio}
        Qual a distância entre esses pontos e a origem
        \begin{alphalist}[3]
            \item $(4,1)$
            \item $(-2,-4)$
            \item $(-5,-2)$
            \item $(1,-10)$
            \item $(6,4)$
            \item $(-3,4)$
            \item $(-1,1)$
            \item $(5,-2)$
        \end{alphalist}
        \begin{answers}
            \begin{mathlist}[4]
                \item \sqrt{17}
                \item 2 \sqrt{5}
                \item \sqrt{29}
                \item \sqrt{101}
                \item 2 \sqrt{13}
                \item 5
                \item \sqrt{2}
                \item \sqrt{29}
            \end{mathlist}
        \end{answers}
    \end{exercicio}
        
    \begin{exercicio}
        Qual a distância de cada ponto ao ponto~$(-1,1)$
        \begin{alphalist}[3]
            \item $(4,1)$
            \item $(-2,-4)$
            \item $(-5,-2)$
            \item $(1,-10)$
            \item $(6,4)$
            \item $(-3,4)$
            \item $(5,-2)$
        \end{alphalist}
        \begin{answers}
            \begin{mathlist}[4]
                \item 5
                \item \sqrt{26}
                \item 5
                \item 5 \sqrt{5}
                \item \sqrt{58}
                \item \sqrt{13}
                \item 3 \sqrt{5}
            \end{mathlist}
        \end{answers}
    \end{exercicio}
    
    \begin{exercicio}
        Determine a distância entre os pontos dados.
        \begin{alphalist}[2]
            \item $(1, 3)$ e $(4,7)$
            \item $(-1,3)$ e $(4,9)$
        \end{alphalist}	
        \begin{answers}
            \begin{mathlist}[2]
                \item 5
                \item \sqrt{61}
            \end{mathlist}
        \end{answers}
    \end{exercicio}
    
    \begin{exercicio}
        Determine as coordenadas dos pontos  que se encontram a 10 unidades da origem e têm coordenada $y$ igual a $-6$.
    \end{exercicio}
    
    \begin{exercicio}
        Determine as coordenadas dos pontos  que se encontram a 5 unidades da origem e têm coordenada $x$ igual a $3$.
    \end{exercicio}
    
    \begin{exercicio}
        Mostre que os pontos $(3, \, 4)$, $(-3, \, 7)$, $(-6, \, 1)$ e $(0, \, -2)$ formam os vértices de um quadrado.
        \dica{Lembre das propriedades dos lados e das diagonais de um quadrado.}
    \end{exercicio}
    
    \begin{exercicio}
        Mostre que o triângulo com vértices $(-5, \, 2)$, $(-2, \, 5)$ e $(5, \, -2)$ é um triângulo retângulo.
        \dica{Lembre da relação que vale para todo triângulo retângulo.}
    \end{exercicio}
        
    \begin{exercicio}
        Calcule a distância entre os pontos $A(a-3,b+4)$ e $B(a+2,b-8)$.
    \end{exercicio}
    
    \begin{exercicio}
        Calcule o perímetro do triângulo $ABC$, sendo dados $A(2,1)$, $B(-1,3)$ e $C(4,-2)$.
    \end{exercicio}
    
    \begin{exercicio}
        Determine $x$ de modo que o triângulo $ABC$ seja retângulo em $B$, sendo dado que $A(4,5)$, $B(1,1)$ e $C(x,4)$.
    \end{exercicio}
    
    \begin{exercicio}
        Se $P(x,y)$ está a mesma distância de $A(-3,7)$ e $B(4,3)$, qual a relação entre $x$ e $y$?
        \begin{answers}
            \tabto{2em} $14x-8y+33=0$
        \end{answers}
    \end{exercicio}

    \begin{exercicio}
        Dados $A(x,5)$, $B(-2,3)$ e $C(4,1)$ obtenha $x$ de modo que $A$ esteja a mesma distância de $B$ e de $C$.
    \end{exercicio}
    
    \begin{exercicio}
        Dados $(-2,4)$ e $(3,-1)$ vértices consecutivos de um quadrado determine os outros dois vértices.
    \end{exercicio}
    
\end{exercicios}

%------------------------------------------------------------------------------%
